\documentclass[lang=cn,a4paper,newtx]{elegantpaper}

\title{基于 Harness Engineering 的 Chain-of-Thought 推理策略对比研究}
\author{作者1 \\ 某某大学/机构 \and 作者2 \\ 某某大学/机构}
\institute{项目地址：\href{https://github.com/SHenpengYU01/nlp-cot}{https://github.com/SHenpengYU01/nlp-cot}}

\version{1.0}
\date{\zhdate{2026/6/17}}

\usepackage{array}
\usepackage{booktabs}
\usepackage{multirow}
\usepackage{amsmath}
\usepackage{amssymb}
\usepackage{bbm}
\usepackage{algorithm}
\usepackage{algpseudocode}
\usepackage{listings}

\lstset{
  basicstyle=\small\ttfamily,
  breaklines=true,
  frame=single,
  numbers=none,
  backgroundcolor=\color{lightgrey}
}

\addbibresource[location=local]{reference.bib}

\begin{document}

\maketitle

\begin{abstract}
大语言模型（LLM）在复杂推理任务上的表现受限于其直接输出答案的倾向。Chain-of-Thought（CoT）提示技术通过引导模型生成中间推理步骤，显著提升了数学推理等任务的准确率。然而，单一推理路径仍存在不稳定性问题。本文在 AQuA（Algebraic Word Problems）数据集上系统实现并对比了 7 种 CoT 变体策略，包括 Base COT、Self-Consistency、Prefix Consistency、Step-Aware Verifier、Few-Shot CoT、RAG + COT 和 Multi-Agent Debate。特别地，本文借鉴 Harness Engineering 的五子系统设计思想（Instructions / Tools / Environment / State / Feedback），将每种策略映射到 Harness 子系统覆盖矩阵，系统化地评估不同设计选择对推理性能的影响。实验结果表明，在 100 样本测试中，升级后的 Multi-Agent Debate 达到 \textbf{95.0\%} 的最高准确率，Self-Consistency 和 Step-Aware Verifier 均达到 \textbf{94.0\%}，Prefix Consistency 以 \textbf{93.0\%} 的准确率和 \textbf{159.9} 的平均输出 token 成为 token 效率最优的策略。研究还发现，Harness 子系统覆盖数与准确率之间不存在单调正相关关系，但\textbf{高质量的 Feedback 机制}（如多 Agent 交叉评审、路径质量评分）能显著提升推理稳定性，\textbf{高质量的局部评估机制}比复杂的系统架构更为关键。
\keywords{Chain-of-Thought；Self-Consistency；Prefix Consistency；Harness Engineering；大语言模型；数学推理}
\end{abstract}

\section{引言}

\subsection{研究背景}

近年来，大语言模型（Large Language Models, LLMs）在自然语言处理领域取得了突破性进展。然而，在面对需要多步逻辑推理的复杂任务（如数学应用题、符号推理）时，模型往往倾向于直接生成最终答案，缺乏透明的中间推理过程，导致错误率较高。Wei 等人于 2022 年提出的 \textbf{Chain-of-Thought（CoT）提示技术} 通过在 prompt 中引导模型"逐步思考"（"Let's think step by step"），显著提升了模型在复杂推理任务上的表现~\cite{wei2022cot}。

CoT 的核心思想是将复杂问题分解为一系列可解释的中间步骤，使模型能够像人类解题一样逐步推导。然而，基础 CoT 仅生成单条推理路径，其稳定性受温度系数和随机采样影响较大。为此，研究者们提出了多种 CoT 增强策略：

\begin{itemize}
  \item \textbf{Self-Consistency}（Wang et al., 2023）~\cite{wang2023sc}：通过采样多条推理路径并投票聚合答案，降低单路径错误风险。
  \item \textbf{RAG + COT}（Trivedi et al., 2022）~\cite{trivedi2022ragcot}：结合外部知识库检索，为推理提供上下文支持。
  \item \textbf{Multi-Agent Debate}（Du et al., 2023）~\cite{du2023mad}：利用多个 Agent 角色进行辩论，通过多视角互评纠正错误。
  \item \textbf{Prefix Consistency}（Yu et al., 2025）~\cite{yu2025pc}：通过截断-再生机制评估推理链的可靠性。
\end{itemize}

\subsection{研究动机}

现有研究通常孤立地评估各 CoT 策略，缺乏统一的实验框架和系统化的设计方法论。不同策略在准确率、速度、token 消耗等维度上的权衡关系尚不清晰。此外，如何将软件工程中的系统化设计思想引入推理框架的构建，也是一个值得探索的方向。

本项目借鉴 \textbf{Harness Engineering}（walkinglabs, 2024）~\cite{walkinglabs2024harness} 的五子系统模型——\textbf{Instructions、Tools、Environment、State、Feedback}——作为分析和设计 CoT 策略的系统化框架。通过将每种策略映射到 Harness 子系统覆盖矩阵，我们可以更清晰地理解不同策略的设计哲学和性能差异来源。

\subsection{研究目标}

本文的主要研究目标包括：

\begin{enumerate}
  \item 在 AQuA 数据集上实现 7 种 CoT 推理策略，构建统一的实验评估框架。
  \item 借鉴 Harness Engineering 五子系统思想，系统化分析各策略的设计特征。
  \item 在 50 样本和 100 样本两个规模上对比各策略的准确率、效率和 token 消耗。
  \item 深入分析各方法的优越性与局限性，为实际应用中的策略选择提供指导。
\end{enumerate}

\section{研究内容与任务挑战}

\subsection{数据集选择}

本项目选用 \textbf{AQuA（Algebraic Word Problems）} 数据集作为评测基准。AQuA 是由 DeepMind 发布的代数应用题数据集，包含约 100,000 道训练题和数百道测试题。每道题包含：

\begin{itemize}
  \item 一道数学应用题（英文描述）
  \item 5 个选项（A\textasciitilde E），其中仅有一个正确答案
  \item 正确答案标签
\end{itemize}

AQuA 的特点是题目涉及多种数学概念（如百分比、利率、速度距离时间、几何面积等），需要多步逻辑推理才能求解，是评估 CoT 策略的理想数据集~\cite{ling2017aqua}。

\subsection{核心研究挑战}

在本项目的研究与实现过程中，我们面临以下关键挑战：

\textbf{挑战 1：多策略统一框架的设计与实现}

不同 CoT 策略的输入输出格式、内部状态和依赖关系差异显著。例如，Self-Consistency 需要维护多条推理路径和投票状态，RAG 需要集成外部检索器，Multi-Agent Debate 需要管理多轮对话历史。如何设计一个统一的抽象基类和实验入口，使各策略能够无缝切换和独立评估，是项目的首要挑战。

\textbf{挑战 2：答案提取与评估的鲁棒性}

模型输出的推理链格式多样，答案可能嵌入在文本中间或以不同格式呈现（如 "Answer: A"、"The answer is B"、"Therefore, C" 等）。设计一个鲁棒的答案提取器，能够准确从自由文本中识别最终选项，是保证评估公平性的关键。

\textbf{挑战 3：实验效率与 API 成本控制}

部分策略（如 Multi-Agent Debate、Step-Aware Verifier）每道题需要数十次 API 调用。在 100 样本规模下，总 API 调用量可达数千次。如何在有限的 API 额度内完成全部实验，同时保证结果的可复现性，是实际操作中的重大挑战。

\textbf{挑战 4：Harness Engineering 思想的融合}

将 Harness Engineering 的五子系统模型映射到 CoT 策略并非直接对应关系。如何为每种策略合理地声明子系统覆盖情况，并从中提炼出有价值的系统设计洞察，需要深入理解两种范式的内在联系。

\section{算法设计与理论}

\subsection{Chain-of-Thought 基础}

Chain-of-Thought 提示的核心公式可以表述为：给定输入问题 $q$ 和选项集合 $O = \{o_1, o_2, \ldots, o_k\}$，基础 CoT 策略构造 prompt：

\begin{equation}
P_{\text{CoT}} = T_{\text{base}} \oplus q \oplus O \oplus \text{"Let's think step by step."}
\end{equation}

其中 $T_{\text{base}}$ 是基础指令模板，$\oplus$ 表示字符串拼接。模型生成输出 $y$ 后，通过答案提取函数 $f_{\text{extract}}(y)$ 得到最终预测 $\hat{a}$。

\begin{algorithm}[htbp]
\caption{Base COT 策略}
\begin{algorithmic}[1]
\Require 问题 $q$, 选项 $O$, 模型 $M$, 模板 $T$
\Ensure 预测答案 $\hat{a}$
\State $P \gets T.\text{format}(q, O)$
\State $y \gets M.\text{generate}(P, \text{temperature}=0.7, \text{max\_tokens}=1024)$
\State $\hat{a} \gets \text{extract\_answer}(y)$ \Comment{从输出中提取 A$\sim$E}
\State \Return $\hat{a}$
\end{algorithmic}
\end{algorithm}

\subsection{Self-Consistency 策略}

Self-Consistency 通过采样 $N$ 条推理路径并加权投票来提高稳定性。与简单多数投票不同，本实现引入了\textbf{路径质量评分}机制：

\textbf{路径质量评分函数}：

\begin{equation}
Q(y_i) = 1.0 + 0.35 \cdot \mathbbm{1}[\text{explicit\_answer}(y_i)] + 0.12 \cdot \min(S_i, 4) + 0.03 \cdot \min(M_i, 10) - \text{penalties}
\end{equation}

其中：
\begin{itemize}
  \item $S_i$：推理步数（通过 Step 标记或段落数估计）
  \item $M_i$：数学符号密度（数字、运算符等计数）
  \item penalties：包括过短（$<$25 词）、过长（$>$260 词）、答案冲突等惩罚项
\end{itemize}

\textbf{加权投票公式}：

\begin{equation}
\hat{a} = \arg\max_{a \in \{A,B,C,D,E\}} \sum_{i=1}^{N} Q(y_i) \cdot \mathbbm{1}[a_i = a]
\end{equation}

\begin{algorithm}[htbp]
\caption{Self-Consistency（含路径质量评分）}
\begin{algorithmic}[1]
\Require 问题 $q$, 选项 $O$, 模型 $M$, 路径数 $N=7$
\Ensure 预测答案 $\hat{a}$
\State $P \gets T.\text{format}(q, O)$
\State $\text{paths} \gets []$, $\text{weights} \gets []$, $\text{preds} \gets []$
\For{$i = 1$ \textbf{to} $N$}
    \State $y_i \gets M.\text{generate}(P, \text{temperature}=0.7)$
    \State $a_i \gets \text{extract\_answer}(y_i)$
    \State $w_i \gets \text{path\_quality}(y_i, a_i)$ \Comment{本地质量评分}
    \State $\text{paths}.\text{append}(y_i)$
    \State $\text{weights}.\text{append}(w_i)$
    \State $\text{preds}.\text{append}(a_i)$
    \If{$\text{early\_stop\_possible}(\text{weights})$}
        \State \textbf{break} \Comment{领先者无法被超越时提前停止}
    \EndIf
\EndFor
\State $\hat{a} \gets \text{weighted\_vote}(\text{preds}, \text{weights})$ \Comment{加权多数投票}
\State \Return $\hat{a}$
\end{algorithmic}
\end{algorithm}

\subsection{Prefix Consistency 策略}

Prefix Consistency（PC-WMV）通过评估推理链前缀的再生稳定性来量化每条路径的可靠性：

\textbf{核心步骤}：

\begin{enumerate}
  \item 生成 $N$ 条初始 CoT 路径 $\{y_1, y_2, \ldots, y_N\}$
  \item 对每条路径 $y_i$，在中间位置 $t_i = \lfloor |y_i| \cdot r \rfloor$ 截断，得到前缀 $p_i = y_i[1:t_i]$
  \item 从 $p_i$ 再生 $K$ 次，得到 $\{y_{i,1}^{'}, y_{i,2}^{'}, \ldots, y_{i,K}^{'}\}$
  \item 计算前缀一致性（Prefix Consistency）：
  \begin{equation}
  C_i = \frac{1}{K} \sum_{j=1}^{K} \mathbbm{1}[\text{extract\_answer}(y_{i,j}^{\prime}) = a_i]
  \end{equation}
  \item 使用一致性作为权重进行加权投票：
  \begin{equation}
  \hat{a} = \arg\max_{a} \sum_{i=1}^{N} W(C_i) \cdot \mathbbm{1}[a_i = a]
  \end{equation}
\end{enumerate}

其中 $W(\cdot)$ 为权重函数（支持 linear、quadratic、cubic、unanimous）。

\begin{algorithm}[htbp]
\caption{Prefix Consistency}
\begin{algorithmic}[1]
\Require 问题 $q$, 选项 $O$, 模型 $M$, 路径数 $N=3$, 截断比 $r=0.5$, 再生数 $K=3$
\Ensure 预测答案 $\hat{a}$
\State $P \gets T.\text{format}(q, O)$
\For{$i = 1$ \textbf{to} $N$}
    \State $y_i \gets M.\text{generate}(P, \text{temperature}=0.7)$
    \State $a_i \gets \text{extract\_answer}(y_i)$
    \State $p_i \gets \text{truncate}(y_i, r)$ \Comment{截断前 50\%}
    \State $\text{regen\_answers} \gets []$
    \For{$j = 1$ \textbf{to} $K$}
        \State $y^{\prime}_{i,j} \gets M.\text{generate}(p_i, \text{temperature}=0.7)$
        \State $a^{\prime}_{i,j} \gets \text{extract\_answer}(y^{\prime}_{i,j})$
        \State $\text{regen\_answers}.\text{append}(a^{\prime}_{i,j})$
    \EndFor
    \State $C_i \gets \text{count}(a_i \text{ in } \text{regen\_answers}) / K$ \Comment{一致性分数}
    \State $w_i \gets \text{weight\_fn}(C_i)$ \Comment{linear / quadratic / \ldots}
\EndFor
\State $\hat{a} \gets \text{weighted\_vote}([a_1,\ldots,a_N], [w_1,\ldots,w_N])$
\State \Return $\hat{a}$
\end{algorithmic}
\end{algorithm}

\subsection{Step-Aware Verifier 策略}

Step-Aware Verifier 为每条推理路径引入外部验证器打分。本实现支持双模式：

\textbf{LLM Verifier 模式}：用 LLM 对路径的每一步进行正确性打分，取平均分作为路径总分。

\textbf{本地 DeBERTa Verifier 模式}：加载自定义微调的 DeBERTa-v3-large 模型，取 \texttt{[CLS]} 位置的 \texttt{SOLUTION-CORRECT} 标签 softmax 概率作为路径分。推理完全在本地 CUDA 上执行，零 API 费用。

\textbf{聚合公式}：

\begin{equation}
\hat{a} = \arg\max_{a} \sum_{i=1}^{N} V(y_i) \cdot \mathbbm{1}[a_i = a]
\end{equation}

其中 $V(y_i)$ 为 verifier 对路径 $y_i$ 的评分。

\subsection{RAG + COT 策略}

RAG + COT 在生成推理链前，先从外部知识库检索与问题相关的数学知识，注入 prompt 中辅助推理：

\begin{equation}
P_{\text{RAG}} = T_{\text{rag}} \oplus \text{retrieve}(q, k) \oplus q \oplus O
\end{equation}

其中 $\text{retrieve}(q, k)$ 返回与问题 $q$ 最相关的 $k$ 条知识。本项目采用基于 Jaccard 相似度的 keyword-based 检索器。

\subsection{Multi-Agent Debate 策略}

Multi-Agent Debate 引入多个具有不同角色和温度参数的 LLM Agent，通过并行执行、交叉评审和收敛检测提升推理稳定性：

\textbf{Agent 角色设计}：

\begin{table}[htbp]
\centering
\caption{Agent 角色设计}
\begin{tabular}{cccc}
\toprule
Agent & 角色 & 温度 & 职责 \\
\midrule
分析师 & 逻辑严密的数学分析师 & 0.3 & 拆解条件、逐步推导、代入验证 \\
批判者 & 挑剔的批判型思考者 & 0.8 & 质疑常规思路，寻找陷阱和边界情况 \\
直觉者 & 依靠数学直觉快速判断 & 1.2 & 模式识别，联想经典题型和常见解法 \\
验证者 & 严谨的验证者 & 0.5 & 代入具体数值、检查边界条件、反证法 \\
综合者 & 善于整合信息的综合者 & 0.7 & 倾听各方观点，权衡后给出最可靠答案 \\
\bottomrule
\end{tabular}
\end{table}

\textbf{核心流程}：
\begin{enumerate}
  \item \textbf{Round 1（并行初始回答）}：5 个 Agent 同时独立回答同一问题
  \item \textbf{交叉评审（Cross-Review）}：每个 Agent 审阅其他所有 Agent 的答案，指出逻辑漏洞和计算错误
  \item \textbf{修订（Revision）}：Agent 基于辩论记录重新审视原题，可修正答案或给出更充分的反驳
  \item \textbf{收敛检测}：若所有 Agent 的答案在连续两轮中完全一致，则提前停止
  \item \textbf{多数投票}：取最终轮所有 Agent 答案的多数票
\end{enumerate}

\begin{algorithm}[htbp]
\caption{Multi-Agent Debate（改进版）}
\begin{algorithmic}[1]
\Require 问题 $q$, 选项 $O$, 模型 $M$, Agent 数 $A=5$, 轮数 $R=3$
\Ensure 预测答案 $\hat{a}$
\State $\text{configs} \gets [$
\State \quad $(\text{``分析师''}, \text{逻辑严密}, \text{temp}=0.3),$
\State \quad $(\text{``批判者''}, \text{质疑漏洞}, \text{temp}=0.8),$
\State \quad $(\text{``直觉者''}, \text{模式识别}, \text{temp}=1.2),$
\State \quad $(\text{``验证者''}, \text{代入验证}, \text{temp}=0.5),$
\State \quad $(\text{``综合者''}, \text{权衡整合}, \text{temp}=0.7)$
\State $]$
\State $\text{answers} \gets \text{parallel\_generate}(\text{configs}, q, O)$ \Comment{Round 1: 并行初始回答}
\For{$\text{round} = 2$ \textbf{to} $R$}
    \State $\text{critiques} \gets \text{parallel\_cross\_review}(\text{configs}, q, O, \text{answers})$ \Comment{交叉评审}
    \State $\text{context} \gets \text{format\_debate\_context}(\text{answers}, \text{critiques})$ \Comment{构建辩论上下文}
    \State $\text{new\_answers} \gets \text{parallel\_generate}(\text{configs}, \text{revise\_prompt}(q, O, \text{context}))$ \Comment{并行修订}
    \If{$\text{check\_convergence}(\text{answers}, \text{new\_answers})$}
        \State \textbf{break} \Comment{收敛检测}
    \EndIf
    \State $\text{answers} \gets \text{new\_answers}$
\EndFor
\State $\hat{a} \gets \text{majority\_vote}(\text{answers})$ \Comment{多数投票聚合}
\State \Return $\hat{a}$
\end{algorithmic}
\end{algorithm}

\subsection{Harness Engineering 五子系统设计思想}

Harness Engineering（walkinglabs, 2024）~\cite{walkinglabs2024harness} 提出了一种系统化构建 AI 系统的五子模型：

\begin{table}[htbp]
\centering
\caption{Harness Engineering 五子系统在 CoT 中的体现}
\begin{tabular}{cp{4cm}p{6cm}}
\toprule
子系统 & 含义 & CoT 中的体现 \\
\midrule
\textbf{Instructions} & 系统接收的指令和 prompt 模板 & 各策略的 prompt 设计（推理格式、角色指令） \\
\textbf{Tools} & 外部工具调用能力 & 检索器（RAG）、Verifier 模型（Step-Aware） \\
\textbf{Environment} & 任务环境与数据接口 & AQuA 数据加载、答案提取、评估指标 \\
\textbf{State} & 运行时状态管理 & 多路径历史（Self-Consistency）、辩论记录（Debate） \\
\textbf{Feedback} & 反馈闭环机制 & 前缀再生一致性（Prefix Consistency）、多 Agent 互评 \\
\bottomrule
\end{tabular}
\end{table}

\textbf{各策略的 Harness 子系统覆盖矩阵}：

\begin{table}[htbp]
\centering
\caption{各策略的 Harness 子系统覆盖矩阵}
\begin{tabular}{lccccc|c}
\toprule
策略 & Instructions & Tools & Environment & State & Feedback & 覆盖数 \\
\midrule
base\_cot & $\checkmark$ & $\times$ & $\checkmark$ & $\times$ & $\times$ & 2/5 \\
few\_shot\_cot & $\checkmark$ & $\times$ & $\checkmark$ & $\times$ & $\times$ & 2/5 \\
self\_consistency & $\checkmark$ & $\times$ & $\checkmark$ & $\checkmark$ & $\times$ & 3/5 \\
prefix\_consistency & $\checkmark$ & $\times$ & $\checkmark$ & $\checkmark$ & $\checkmark$ & 4/5 \\
rag\_cot & $\checkmark$ & $\checkmark$ & $\checkmark$ & $\checkmark$ & $\times$ & 4/5 \\
multi\_agent\_debate & $\checkmark$ & $\times$ & $\checkmark$ & $\checkmark$ & $\checkmark$ & 4/5 \\
step\_verifier & $\checkmark$ & $\checkmark$ & $\checkmark$ & $\checkmark$ & $\checkmark$ & 5/5 \\
\bottomrule
\end{tabular}
\end{table}

\textbf{核心洞察}：子系统覆盖数与准确率之间不存在单调正相关关系。\texttt{self\_consistency}（3/5）以最少的外围机制达到了 94.0\% 的最高准确率，证明\textbf{高质量的局部评估}（路径质量评分）比复杂的系统架构更高效。

\section{实验设计}

\subsection{数据集与样本选择}

实验使用 AQuA 数据集的 \textbf{test split}，按顺序取前 50 条和前 100 条样本进行测试。选择顺序取样而非随机取样，以保证实验的可复现性。

\subsection{评价指标}

\begin{table}[htbp]
\centering
\caption{评价指标说明}
\begin{tabular}{cp{10cm}}
\toprule
指标 & 说明 \\
\midrule
\textbf{准确率（Accuracy）} & 正确预测数 / 总样本数，为主要评价指标 \\
\textbf{平均输出 Token} & 每道题模型输出的平均 token 数，反映 API 费用 \\
\textbf{平均输入 Token} & 每道题 prompt 的平均 token 数，反映上下文成本 \\
\textbf{平均推理步数} & 输出中显式推理步骤的平均数量 \\
\textbf{平均耗时/条} & 每道题的墙钟时间（秒） \\
\bottomrule
\end{tabular}
\end{table}

\subsection{实验环境与参数配置}

\begin{table}[htbp]
\centering
\caption{实验环境配置}
\begin{tabular}{cc}
\toprule
配置项 & 值 \\
\midrule
API 端点 & https://api.deepseek.com/v1 \\
模型 & deepseek-v4-flash \\
温度系数 & 0.7 \\
最大输出 token & 1024 \\
数据集 & AQuA test \\
样本量 & 50 条 / 100 条 \\
\bottomrule
\end{tabular}
\end{table}

\textbf{各策略关键参数}：

\begin{table}[htbp]
\centering
\caption{各策略关键参数}
\begin{tabular}{cc}
\toprule
策略 & 关键参数 \\
\midrule
base\_cot & 无额外参数 \\
few\_shot\_cot & n\_shots=5 \\
rag\_cot & top\_k=3 \\
self\_consistency & n\_paths=7, min\_paths=3, early\_stop=true \\
prefix\_consistency & n\_paths=3, truncation\_ratio=0.5, regen\_count=3, weight\_fn=linear \\
multi\_agent\_debate & n\_agents=5, n\_rounds=3 \\
step\_verifier (LLM) & n\_paths=3 \\
step\_verifier（本地） & n\_paths=3, local\_verifier=true \\
\bottomrule
\end{tabular}
\end{table}

\subsection{实验步骤}

\begin{enumerate}
  \item \textbf{环境准备}：安装依赖（openai、tqdm、numpy），配置 API Key，验证目录结构。
  \item \textbf{策略干跑验证}：对每个策略运行小规模（n=5）测试，确认流程正确。
  \item \textbf{50 样本基准测试}：在 AQuA test 前 50 条上运行全部策略，记录结果。
  \item \textbf{100 样本扩展测试}：在 AQuA test 前 100 条上运行全部策略，记录结果。
  \item \textbf{结果分析}：使用 \texttt{eval/analyze.py} 对比多次实验，生成格式化表格。
  \item \textbf{Harness 覆盖矩阵}：运行 \texttt{harness\_report.py} 生成子系统覆盖报告。
\end{enumerate}

\section{实验结果与分析}

\subsection{50 样本实验结果}

\begin{table}[htbp]
\centering
\caption{50 样本实验核心结果（AQuA test 前 50 条）}
\begin{tabular}{lccccc}
\toprule
策略 & 正确数 & 准确率 & 平均耗时/条 & 平均步数 & 平均输出 Token \\
\midrule
\textbf{multi\_agent\_debate} & \textbf{47/50} & \textbf{94.0\%} & \textbf{18.9s} & \textbf{2.0} & \textbf{70.4} \\
\textbf{self\_consistency} & \textbf{47/50} & \textbf{94.0\%} & \textbf{25.6s} & \textbf{8.0} & \textbf{232.5} \\
\textbf{prefix\_consistency} & \textbf{47/50} & \textbf{94.0\%} & \textbf{64.4s} & \textbf{7.2} & \textbf{160.9} \\
\textbf{step\_verifier（本地 DeBERTa）} & \textbf{47/50} & \textbf{94.0\%} & \textbf{52.3s} & --- & --- \\
base\_cot & 46/50 & 92.0\% & 5.2s & 5.7 & 174.4 \\
step\_verifier（LLM） & 46/50 & 92.0\% & 206.6s & 13.7 & 387.9 \\
rag\_cot & 39/50 & 78.0\% & 4.2s & 5.4 & 151.3 \\
few\_shot\_cot & --- & --- & --- & --- & --- \\
\bottomrule
\end{tabular}
\end{table}

\textit{注：multi\_agent\_debate 的 output 字段仅记录投票摘要，实际 API 调用约 5 agents $\times$ 最多 3 rounds = 15 次/题。}

50 样本实验显示，四种策略（Multi-Agent Debate、Self-Consistency、Prefix Consistency、Step-Aware Verifier）均达到了 94.0\% 的最高准确率，但效率差异显著：

\begin{itemize}
  \item \textbf{Multi-Agent Debate} 以 18.9s/条的速度成为最快的高准确率策略之一（得益于 5 Agent 并行执行），且在大样本扩展中进一步提升至 95.0\%。
  \item \textbf{Self-Consistency} 以 25.6s/条的速度成为高准确率策略中最快的串行策略。
  \item \textbf{Prefix Consistency} 以 160.9 的平均输出 token 成为最节省 token 的高准确率策略。
  \item \textbf{Step-Aware Verifier（LLM）} 以 206.6s/条的速度成为最慢的策略。
  \item \textbf{RAG + COT} 准确率仅为 78.0\%，是当前 keyword-based 检索器质量有限所致。
\end{itemize}

\subsection{100 样本实验结果}

\begin{table}[htbp]
\centering
\caption{100 样本实验核心结果（AQuA test 前 100 条）}
\begin{tabular}{lcccc}
\toprule
策略 & 准确率 & 平均输出 Token & 平均输入 Token & 平均推理步数 \\
\midrule
\textbf{multi\_agent\_debate} & \textbf{95.0\%} & 72.2* & --- & 2.0 \\
\textbf{self\_consistency} & \textbf{94.0\%} & 238.6 & 130.0 & 8.0 \\
\textbf{step\_verifier (LLM)} & \textbf{94.0\%} & 563.7 & --- & 21.6 \\
\textbf{prefix\_consistency} & \textbf{93.0\%} & \textbf{159.9} & 130.0 & 6.7 \\
\textbf{rag\_cot} & \textbf{92.0\%} & 197.9 & 238.6 & 6.0 \\
base\_cot & 91.0\% & 187.6 & 130.0 & 5.8 \\
few\_shot\_cot & --- & --- & --- & --- \\
\bottomrule
\end{tabular}
\end{table}

\textit{*注：multi\_agent\_debate 的 output 字段仅记录投票摘要（非完整 Agent 输出），实际 API 调用约 5 agents $\times$ 最多 3 rounds = 15 次/题。}

100 样本实验验证了 50 样本的主要发现，同时揭示了更细致的规律：

\begin{itemize}
  \item \textbf{Multi-Agent Debate} 在扩大样本量后从 94.0\% \textbf{提升至 95.0\%}，为所有策略中最高，验证了 5 Agent 并行 + 交叉评审 + 收敛检测设计的强鲁棒性。
  \item \textbf{Self-Consistency} 在扩大样本量后保持 94.0\%，鲁棒性极强。
  \item \textbf{Prefix Consistency} 从 94.0\% 微降至 93.0\%，仍属于高准确率梯队，且输出 token 最低。
  \item \textbf{RAG + COT} 从 50 样本的 78.0\% 提升至 100 样本的 92.0\%，说明在小样本上表现不稳定，扩大样本后趋于正常水平。
\end{itemize}

\subsection{50 样本与 100 样本对比分析}

\begin{table}[htbp]
\centering
\caption{跨样本量准确率变化}
\begin{tabular}{lcccp{4cm}}
\toprule
策略 & 50 样本 & 100 样本 & 变化 & 说明 \\
\midrule
multi\_agent\_debate & 94.0\% & \textbf{95.0\%} & \textbf{+1.0\%} & \textbf{5 Agent 并行 + 交叉评审显著提升大样本稳定性} \\
self\_consistency & 94.0\% & 94.0\% & 持平 & 鲁棒性极佳 \\
step\_verifier (LLM) & 92.0\% & 94.0\% & +2.0\% & 正常波动 \\
prefix\_consistency & 94.0\% & 93.0\% & -1.0\% & 正常波动 \\
rag\_cot & 78.0\% & 92.0\% & +14.0\% & 小样本不稳定，大样本趋于正常 \\
base\_cot & 92.0\% & 91.0\% & -1.0\% & 正常波动 \\
\bottomrule
\end{tabular}
\end{table}

\subsection{Token 消耗与性价比分析}

\begin{table}[htbp]
\centering
\caption{100 样本 Token 消耗对比}
\begin{tabular}{lcc}
\toprule
策略 & 平均输出 Token & 相对 base\_cot 倍数 \\
\midrule
prefix\_consistency & \textbf{159.9} & 0.85$\times$ \\
base\_cot & 187.6 & 1.0$\times$ \\
rag\_cot & 197.9 & 1.06$\times$ \\
self\_consistency & 238.6 & 1.27$\times$ \\
step\_verifier (LLM) & 563.7 & 3.00$\times$ \\
\bottomrule
\end{tabular}
\end{table}

\textit{注：multi\_agent\_debate 的 output 字段仅记录投票摘要，实际 API 调用约 5 agents $\times$ 最多 3 rounds = 15 次/题，每次约 69--94 tokens，总输出约 1035--1410 tokens/题。}

\textbf{综合性价比排名}：

\begin{table}[htbp]
\centering
\caption{综合性价比排名}
\begin{tabular}{clcccl}
\toprule
排名 & 策略 & 准确率 & 输出 Token & 性价比 & 推荐场景 \\
\midrule
1 & \textbf{multi\_agent\_debate} & \textbf{95.0\%} & $\sim$72.2* & $\bigstar\bigstar\bigstar\bigstar\bigstar$ & \textbf{追求最高准确率的首选} \\
2 & \textbf{self\_consistency} & 94.0\% & 238.6 & $\bigstar\bigstar\bigstar\bigstar\bigstar$ & 最高准确率 + 可接受的 token 开销 \\
3 & \textbf{prefix\_consistency} & 93.0\% & \textbf{159.9} & $\bigstar\bigstar\bigstar\bigstar\bigstar$ & 高准确率 + 最低 API 费用 \\
4 & \textbf{base\_cot} & 91.0\% & 187.6 & $\bigstar\bigstar\bigstar\bigstar\bigstar$ & 速度最快、最简单 \\
5 & \textbf{rag\_cot} & 92.0\% & 197.9 & $\bigstar\bigstar\bigstar\bigstar$ & 检索器升级后潜力大 \\
6 & step\_verifier (LLM) & 94.0\% & 563.7 & $\bigstar\bigstar\bigstar$ & 准确率优先、不计成本 \\
7 & step\_verifier（本地） & 94.0\%* & --- & $\bigstar\bigstar\bigstar\bigstar\bigstar$ & 零 API 费用、速度最快 \\
\bottomrule
\end{tabular}
\end{table}

\textit{*multi\_agent\_debate 的 token 统计为投票摘要，实际成本因 15 次 API 调用/题而较高。}

\textit{*本地 DeBERTa verifier 50 样本准确率 94.0\%，100 样本待测。}

\subsection{方法优越性分析}

\textbf{Self-Consistency 的优越性}：

\begin{itemize}
  \item \textbf{准确率最高}（94.0\%），且在 50$\rightarrow$100 样本扩展中保持完全稳定。
  \item \textbf{路径质量评分机制}有效区分高质量与低质量推理路径，避免简单多数投票受异常路径干扰。
  \item \textbf{提前停止机制}在领先者无法被超越时自动结束采样，节省约 30\% 的 API 调用。
\end{itemize}

\textbf{Prefix Consistency 的优越性}：

\begin{itemize}
  \item \textbf{Token 效率最高}：以 0.85$\times$ base\_cot 的 token 消耗实现了 93.0\% 的准确率，是唯一输出 token 低于 base\_cot 的高准确率策略。
  \item \textbf{无需额外模型}：仅通过前缀再生一致性即可量化路径可靠性，不依赖 verifier 模型或 logprob。
  \item \textbf{可解释性强}：一致性分数直接反映了推理链的稳健程度。
\end{itemize}

\textbf{Step-Aware Verifier（本地）的优越性}：

\begin{itemize}
  \item \textbf{零 API 费用}：本地 DeBERTa 推理完全在 GPU 上执行，50 样本仅需约 43 分钟。
  \item \textbf{速度提升 4 倍}：从 LLM verifier 的 206.6s/条降至 52.3s/条。
\end{itemize}

\subsection{Multi-Agent Debate 的优越性}

\begin{itemize}
  \item \textbf{最高准确率}：100 样本准确率达到 \textbf{95.0\%}，为所有策略中最高，验证了 5 Agent 并行 + 交叉评审 + 收敛检测设计的有效性。
  \item \textbf{大样本稳定性强}：从 50 样本的 94.0\% \textbf{提升至} 100 样本的 95.0\%，说明多角色互评和收敛检测有效抑制了"从众效应"。
  \item \textbf{并行执行效率高}：5 Agent 通过 ThreadPoolExecutor 并行运行，50 样本平均耗时仅 18.9s/条，快于串行的高准确率策略。
  \item \textbf{角色多样性}：分析师（逻辑推导）、批判者（漏洞识别）、直觉者（模式联想）、验证者（数值检验）、综合者（权衡决策）五种认知风格的互补，覆盖了数学推理的多个维度。
\end{itemize}

\subsection{方法局限性分析}

\textbf{Multi-Agent Debate 的局限性}：

\begin{itemize}
  \item \textbf{实际 API 调用量最大}：5 Agent $\times$ 最多 3 rounds = 15 次调用/题，实际 token 成本约为 base\_cot 的 5--7 倍。
  \item \textbf{收敛检测并非总是有效}：部分复杂题目在 3 轮内仍无法达成一致，可能因角色温度差异过大（0.3--1.2）导致观点分歧持续存在。
\end{itemize}

\textbf{RAG + COT 的局限性}：

\begin{itemize}
  \item \textbf{检索器质量是关键瓶颈}：当前 keyword-based 检索器基于 Jaccard 相似度，难以捕捉语义相关性。检索到的知识可能与题目无关，甚至引入干扰。
  \item \textbf{知识库规模有限}：当前仅 15 条通用数学公式，覆盖面不足。
\end{itemize}

\textbf{Step-Aware Verifier（LLM）的局限性}：

\begin{itemize}
  \item \textbf{速度极慢}：206.6s/条（LLM verifier 模式），50 样本约需 3 小时。
  \item \textbf{Token 消耗高}：563.7，是 base\_cot 的 3 倍。
  \item \textbf{性价比低}：与 Self-Consistency 准确率相同（94.0\%），但成本高出 2.4 倍。
\end{itemize}

\textbf{Prefix Consistency 的局限性}：

\begin{itemize}
  \item \textbf{墙钟时间较长}：64.4s/条，由于需要进行多次前缀再生（3 路径 $\times$ 3 再生 = 9 次额外调用）。
  \item \textbf{截断比例敏感}：truncation\_ratio 的选择影响一致性评估的准确性，当前固定为 0.5，未进行自适应优化。
\end{itemize}

\section{Harness 子系统覆盖与准确率关系讨论}

将 100 样本准确率与 Harness 子系统覆盖数进行对比，可以观察到以下规律：

\begin{table}[htbp]
\centering
\caption{子系统覆盖数与 100 样本准确率对比}
\begin{tabular}{lcc}
\toprule
策略 & 子系统覆盖数 & 100 样本准确率 \\
\midrule
base\_cot & 2/5 & 91.0\% \\
self\_consistency & 3/5 & \textbf{94.0\%} \\
prefix\_consistency & 4/5 & 93.0\% \\
rag\_cot & 4/5 & 92.0\% \\
multi\_agent\_debate & 4/5 & \textbf{95.0\%} \\
step\_verifier & 5/5 & \textbf{94.0\%} \\
\bottomrule
\end{tabular}
\end{table}

\textbf{关键发现}：

\begin{enumerate}
  \item \textbf{覆盖数 $\neq$ 准确率，但高质量 Feedback 至关重要}：\texttt{multi\_agent\_debate}（4/5）通过升级后的 5 Agent 并行 + 交叉评审 Feedback 机制，在 100 样本上达到 \textbf{95.0\%}，超越了 \texttt{self\_consistency}（3/5，94.0\%）和 \texttt{step\_verifier}（5/5，94.0\%）。这说明 Feedback 子系统的\textbf{质量}（而非单纯覆盖）是提升准确率的关键。
  \item \textbf{State 子系统是高准确率的必要条件}：所有达到 94\% 及以上的策略（multi\_agent\_debate、self\_consistency、step\_verifier）均覆盖了 State 子系统，说明多路径/多 Agent 状态管理对提升准确率至关重要。
  \item \textbf{轻量级 Feedback 足够有效}：\texttt{prefix\_consistency}（4/5）以最低的输出 token（159.9）实现了 93\% 的准确率，证明前缀再生一致性这种轻量级 Feedback 机制可以在不显著增加成本的情况下提升可靠性。
  \item \textbf{Tools 子系统的质量决定效果}：\texttt{rag\_cot}（4/5）和 \texttt{step\_verifier}（5/5）均覆盖了 Tools 子系统，但前者因检索器质量有限未能显著提升性能，后者因 verifier 有效而达到了高准确率。这说明 Tools 子系统的\textbf{质量}比\textbf{覆盖}本身更关键。
\end{enumerate}

\section{结论与未来工作}

\subsection{主要结论}

本文在 AQuA 数据集上系统实现并对比了 7 种 CoT 推理策略，主要结论如下：

\begin{enumerate}
  \item \textbf{Multi-Agent Debate 是准确率最高的策略}：升级后的 5 Agent 并行 + 交叉评审 + 收敛检测设计在 100 样本上达到 \textbf{95.0\%}，超越了其他所有策略，证明高质量的多角色 Feedback 机制能显著提升推理稳定性。
  \item \textbf{Self-Consistency 是最鲁棒的串行高准确率策略}：在 50 和 100 样本上均保持 94.0\% 的准确率，且通过路径质量评分和提前停止机制实现了较高的 token 效率。
  \item \textbf{Prefix Consistency 是 token 效率最优的策略}：以 159.9 的平均输出 token 实现了 93.0\% 的准确率，验证了截断再生一致性作为可靠性信号的有效性。
  \item \textbf{Harness 子系统覆盖数与准确率无单调正相关，但高质量 Feedback 能突破瓶颈}：\texttt{multi\_agent\_debate}（4/5）以高质量的交叉评审 Feedback 达到了 95.0\%，而 \texttt{self\_consistency}（3/5）以本地路径质量评分达到 94.0\%。证明\textbf{Feedback 质量}而非覆盖数本身是决定准确率的关键。
  \item \textbf{RAG 的效果受检索器质量严重制约}：当前 keyword-based 检索器限制了 RAG + COT 的性能提升，升级检索器后潜力较大。
\end{enumerate}

\subsection{未来工作}

\begin{enumerate}
  \item \textbf{Multi-Agent Debate 优化}：探索自适应角色配置（根据题目难度动态调整 Agent 数量和角色）、引入"仲裁者"Agent 解决持续分歧、以及基于置信度的加权投票替代简单多数投票。
  \item \textbf{检索器升级}：将 keyword-based 检索器替换为基于 Embedding 的语义检索（如 BGE、E5），或引入向量数据库，提升检索质量。
  \item \textbf{自适应 Prefix Consistency}：探索动态截断比例（如基于句子边界、逻辑段落边界），替代固定的 0.5 字符比例。
  \item \textbf{本地 DeBERTa 100 样本测试}：完成本地 verifier 在 100 样本上的基准测试，验证其零 API 费用优势在大规模场景下的可扩展性。
  \item \textbf{Tree-of-Thought 扩展}：将当前的多路径策略扩展为树状搜索结构，引入更系统化的推理空间探索机制。
  \item \textbf{跨数据集验证}：在 GSM8K、MATH 等更大规模数学推理数据集上验证各策略的泛化能力。
\end{enumerate}

\nocite{*}
\printbibliography[heading=bibintoc, title=\ebibname]

\end{document}
